\documentclass[12pt]{article}
\usepackage[T1]{fontenc}
\usepackage{amsmath,amssymb,amsthm}
\usepackage{hyperref}
\usepackage{geometry}
\usepackage{enumerate}
\geometry{margin=1in}

\newtheorem{definition}{Definition}
\newtheorem{proposition}{Proposition}
\newtheorem{corollary}{Corollary}
\newtheorem{conjecture}{Conjecture}
\newtheorem{remark}{Remark}

\title{Abductive Accumulation:\\
Why Interpretive Errors Are Not Eliminated, Only Suppressed}

\author{Minamo Minamoto\\
Independent\\
\texttt{minamominamoto4f5683f6@gmail.com}}

\date{April 2026}

\begin{document}

\maketitle

\begin{abstract}
The correction model of abduction holds that erroneous hypotheses
are detected and replaced.
We argue that this model is structurally inadequate in a precise sense:
prior interpretive frames are not overwritten but \emph{archived}.
When a cognitive agent adopts a new hypothesis, the prior hypothesis
is not eliminated; it is demoted to a latent layer that remains
accessible and may be reactivated under appropriate conditions.
What changes is not the presence of the prior frame but the
\emph{selection rule} that determines which frame is foregrounded.
We call this process \emph{abductive accumulation} and the resulting
representational structure a \emph{layered interpretive stack}.
This account explains why abductive error is not correctable:
the error is not a missing piece but a latent frame competing for
selection.
The practical consequence is that reliable cognition requires not
the ability to delete erroneous frames---which is impossible---but
the ability to suppress their reactivation: a \emph{selection
stability} condition distinct from but related to the negative
abductive capability introduced in prior work
\cite{minamoto2026overwrite}.
We formalize the layered interpretive stack, the selection rule,
and the conditions under which suppression fails, and derive
implications for the understanding of cognitive bias, clinical
reasoning errors, and the design of trustworthy AI systems.
\end{abstract}

\tableofcontents

% ============================================================
\section{Introduction}
% ============================================================

\paragraph{The correction assumption.}
The standard model of abduction---inference to the best explanation
\cite{harman1965}---presents hypothesis revision as correction:
a prior hypothesis $h_t$ is found inadequate, and a new hypothesis
$h_{t+1}$ is installed in its place.
The prior hypothesis, on this picture, is simply gone.
The agent's representational state after revision is characterized
by $h_{t+1}$ alone; $h_t$ plays no further role.

\paragraph{The residue phenomenon.}
This picture conflicts with a well-documented class of phenomena.
Human reasoners who have been corrected on a factual belief
nonetheless revert to the prior belief under time pressure, emotional
arousal, or situations structurally similar to those in which the
prior belief was formed \cite{ecker2022}.
Eyewitness testimony is unreliable not because new information
is unavailable but because prior perceptual frames compete with
later corrective accounts \cite{loftus1979}.
Anchoring effects persist after anchor correction has been explicitly
performed \cite{epley2001}.
In each case, the erroneous prior frame is not absent---it is
\emph{latent and reactivatable}.

\paragraph{The central claim.}
We argue that this class of phenomena reflects a structural feature
of abductive cognition: interpretive frames are not overwritten but
\emph{archived}.
The representational state of a cognitive agent is not a single
active hypothesis but a \emph{layered interpretive stack}---an
ordered collection of frames, each associated with the evidence
conditions under which it was formed.
When a new frame is adopted, the prior frame is not deleted;
it is pushed down the stack.
What changes is the \emph{selection rule}: the function that
determines, given current context, which frame is foregrounded
for use in perception, classification, and inference.

\paragraph{Consequence for the theory of abductive error.}
If prior frames are archived rather than deleted, then abductive
error is not correctable in the sense of the correction model.
A corrective intervention that installs $h_{t+1}$ does not eliminate
$h_t$; it adds $h_{t+1}$ to the stack and modifies the selection
rule.
Under conditions that activate the selection criteria associated
with $h_t$---time pressure, high similarity to original encoding
context, emotional state, attentional load---$h_t$ may be
foregrounded again.
The error persists as a latent selection risk.

\paragraph{Relation to prior work.}
This paper is the fourth in a series.
Paper 1 \cite{minamoto2026topology} formalizes symbol grounding
as a topological problem.
Paper 2 \cite{minamoto2026cgla} establishes hardware implementation
conditions for grounding.
Paper 3 \cite{minamoto2026overwrite} introduces the overwrite model
of abduction and negative abductive capability (NAC).
The present paper refines Paper 3: the overwrite model is a
first approximation in which the prior frame is treated as
eliminated upon transition.
We now show that the prior frame is in fact archived, and that
this archival structure is the structural basis of abductive
non-correctability.

\begin{remark}[Relation to Paper 3: refinement, not refutation]
Paper 3's overwrite model correctly describes the \emph{selection rule
change} that occurs at commitment time $t^*$: a new interpretive frame
$F_{t+1}$ is installed as the default-active frame, and the agent's
subsequent behavior is governed by $F_{t+1}$ rather than $F_t$.
This description remains accurate. The present paper adds a structural
claim about what happens to $F_t$ after the overwrite: rather than
being deleted, $F_t$ is archived in the interpretive stack and remains
selectable under appropriate environmental conditions. Paper 3 provides
the pre-commitment account (NAC as the capacity to withhold premature
overwrite); the present paper provides the post-commitment account (SSC
as the capacity to suppress reactivation of archived frames). The two
papers are therefore complementary descriptions of different phases of
the same process, not competing accounts.
\end{remark}

\medskip
\noindent\textbf{Central thesis.}
\emph{Abductive error is not fully correctable because prior interpretive
frames are not overwritten but archived; what changes is only the
selection rule. Reliable cognition therefore requires not the
ability to delete erroneous frames but the ability to suppress
their reactivation.}

% ============================================================
\section{The Layered Interpretive Stack}
% ============================================================

\subsection{From Single Frame to Stack}

In Paper 3 \cite{minamoto2026overwrite}, the representational state
of an agent was modeled as a single interpretive frame $F(R)$.
An overwrite transition $R_t \to R_{t+1}$ was defined as a state
change in which $F(R_{t+1}) \not\approx F(R_t)$.

We now replace the single-frame model with a richer structure.

\begin{definition}[Interpretive Stack]
An interpretive stack $\mathcal{L} = (F_1, F_2, \ldots, F_n)$ is
an ordered sequence of interpretive frames, where $F_1$ is the
most recently adopted frame (top of stack) and $F_n$ is the
earliest (bottom of stack).
Each frame $F_k$ is associated with an \emph{acquisition context}
$C_k$: the set of evidence conditions, attentional states, and
temporal markers under which $F_k$ was formed.
\end{definition}

\begin{definition}[Selection Rule]
A selection rule $\mathcal{S} : \mathcal{E} \to \{1, \ldots, n\}$
maps the current environmental state $e \in \mathcal{E}$---including
evidence quality, time pressure, emotional arousal, and contextual
similarity to prior acquisition contexts---to an index $k$,
selecting frame $F_k$ as the active interpretive frame for that
moment.
A canonical functional form is:
\[
\mathcal{S}(e) = \operatorname{arg\,max}_k\;
\mathrm{softmax}_k\!\left(\frac{s(e,\, C_k)}{\tau}\right)
\]
where $C_k$ is the acquisition context of $F_k$,
$s(\cdot,\cdot): \mathcal{E} \times \mathcal{E} \to [0,1]$ is a context-similarity function,
and $\tau > 0$ is a temperature parameter controlling selection sharpness.
High $\tau$ produces near-uniform selection (reactivation risk increases);
low $\tau$ produces sharp selection of the contextually best-matching frame.
\end{definition}

\begin{remark}[The temperature parameter $\tau$ and cognitive load]
The temperature parameter $\tau$ has a concrete cognitive interpretation.
Under normal operating conditions---attentive processing, low time pressure,
familiar but non-overlapping context---$\tau$ is small and selection is sharp:
$F_1$ (the top-of-stack corrected frame) dominates with high probability.
As cognitive load increases, $\tau$ rises and the selection distribution
flattens, increasing the probability that archived frames $F_k$ ($k > 1$)
are activated.
Concretely, in LLMs: very long contexts require integrating over more
tokens, effectively raising $\tau$ for context-similarity estimation;
high-perplexity inputs lie far from the fine-tuning distribution,
degrading $\delta(\cdot, C_1)$ and shifting probability toward earlier
acquisition contexts; removal of chain-of-thought scaffolding removes
the structured cue that normally suppresses $F_k$ activation.
In biological agents: time pressure, emotional arousal, and divided
attention correspond directly to elevated $\tau$ (degraded selection
rule application), producing the load-dependent reactivation effect
documented in the continued influence literature \cite{ecker2022}.
This mapping from $\tau$ to observable cognitive conditions makes
Prediction~3 (Load-Dependent Reactivation) directly testable:
any experimental manipulation that increases $\tau$ should increase
archived-frame reactivation rate.
\end{remark}


\begin{remark}
The correction model corresponds to the degenerate case in which
the stack has depth 1 at all times: each overwrite replaces the
stack with a singleton containing only $F_{t+1}$.
We deny this degeneracy.
\end{remark}

\subsection{Abductive Accumulation}

\begin{definition}[Abductive Accumulation]
An abductive transition is an \emph{accumulation} rather than an
overwrite if the prior frame $F_t$ is archived (pushed onto the
stack) rather than deleted when the new frame $F_{t+1}$ is adopted.
The stack after accumulation is $(F_{t+1}, F_t, \ldots)$.
\end{definition}

\begin{proposition}[Non-Correctability and Reactivation Bound]
\label{prop:non-correct}
Let $\mathcal{L} = (F_{t+1}, F_t, \ldots)$ be the interpretive stack
after a corrective accumulation at time $t+1$.
Then:
\begin{enumerate}[(i)]
    \item $F_t$ remains present in $\mathcal{L}$ and can be selected
    by $\mathcal{S}$ under environmental conditions $e$ such that
    $\mathcal{S}(e) = k$ where $F_k = F_t$.
    The ``correction'' does not eliminate $F_t$; it installs $F_{t+1}$
    as the default-selected frame and modifies the selection rule.
    \item The reactivation probability of $F_t$ is bounded away from
    zero for all inputs $e$ (since $\mathrm{softmax}(\cdot) > 0$
    always), and is highest when $s(e, C_t)$ is large---that is,
    when $e$ closely resembles the acquisition context $C_t$ of $F_t$.
    Formally, $P(\mathcal{S}(e) = k \mid F_k = F_t) > 0$ for all $e$,
    with $P$ increasing monotonically in $s(e, C_t)$.
    This bound cannot be driven to zero by modifying $F_{t+1}$
    alone, since $\mathcal{S}$ responds to $s(e, C_t)$,
    not to the content of $F_{t+1}$.
\end{enumerate}
Under conditions that activate $F_t$'s acquisition context $C_t$,
$F_t$ may be foregrounded, reinstating the prior bias. The agent's
ability to prevent this reactivation defines Selection Stability
Capacity (SSC), introduced in Section~3.
\end{proposition}

\begin{corollary}
Abductive error is not a deleted state.
It is a latent frame with a non-zero selection probability under
a subset of environmental conditions.
``Correcting'' an abductive error reduces but does not eliminate
this selection probability.
\end{corollary}

% ============================================================
\section{Selection Instability as the Source of Bias}
% ============================================================

\begin{definition}[Selection Stability]
A cognitive agent exhibits \emph{selection stability} with respect
to hypothesis $h^*$ if the selection rule $\mathcal{S}$ assigns
$h^*$'s frame to the top index across all relevant environmental
conditions $e \in \mathcal{E}$:
$\mathcal{S}(e) = 1$ for all $e \in \mathcal{E}$.
\end{definition}

\begin{proposition}[Bias as Selection Instability]
\label{prop:bias-selection}
Cognitive bias is not the presence of an erroneous frame in the
stack but the failure of selection stability: the erroneous frame
$F_\mathrm{err}$ is selected under a non-negligible subset of
environmental conditions, producing systematic distortion of
perception, classification, or inference under those conditions.
\end{proposition}

\begin{remark}
This reformulation has a direct implication for debiasing
interventions.
Interventions that provide corrective information---installing
a new frame $F^*$ at the top of the stack---are insufficient
if they do not also modify the selection rule to suppress
$F_\mathrm{err}$ reactivation.
The corrective information may be present in the stack and yet
fail to be selected under the conditions that matter most.
\end{remark}

% ============================================================
\section{Relation to Negative Abductive Capability}
% ============================================================

In Paper 3 \cite{minamoto2026overwrite}, negative abductive
capability (NAC) was defined as the capacity to maintain a
suspension state---holding competing hypotheses without committing
to any single frame---until evidence warrants commitment.

The present account adds a new dimension.
NAC, as defined in Paper 3, addresses the \emph{pre-commitment}
phase: it prevents premature stacking.
The present paper identifies a complementary \emph{post-commitment}
capacity: the ability to suppress reactivation of archived frames
after commitment has occurred.

\begin{definition}[Selection Stability Capacity (SSC)]
An agent possesses \emph{selection stability capacity} if it can
maintain $\mathcal{S}(e) = 1$ (top-of-stack selection) across
the range of environmental conditions $\mathcal{E}$ relevant to
a target task, even when the stack contains archived frames
with overlapping acquisition contexts.
\end{definition}

\begin{remark}
NAC and SSC are distinct but related capacities.
NAC prevents premature accumulation of frames with high error risk.
SSC prevents reactivation of accumulated frames with high error risk.
Together they constitute the structural conditions for reliable
abductive cognition:
\begin{enumerate}
  \item[\textbf{(NAC)}] Suspend commitment until evidence warrants
  stacking.
  \item[\textbf{(SSC)}] Once committed, suppress reactivation of
  archived error frames under conditions that would otherwise
  trigger them.
\end{enumerate}
\end{remark}

% ============================================================
\section{Evidence and Illustrations}
% ============================================================

\subsection{The Continued Influence Effect}

The continued influence effect \cite{ecker2022} refers to the
persistence of misinformation in reasoning even after explicit
correction.
Subjects who are told that a belief was false continue to reason
as if it were true, particularly under cognitive load.

On the accumulation model: the corrective information installs a
new frame $F^*$ at the top of the stack, but the prior frame
$F_\mathrm{err}$ is archived with acquisition context $C_\mathrm{err}$.
Under cognitive load---which reduces the computational resources
available for selection rule application---the selection process
defaults to simpler criteria, reactivating $F_\mathrm{err}$.

\subsection{Anchoring and Adjustment}

Anchoring effects \cite{epley2001} arise when an initial numerical
estimate biases subsequent judgment even after adjustment information
is provided.

On the accumulation model: the anchor establishes $F_\mathrm{anchor}$
as the initial frame. Adjustment installs $F_\mathrm{adj}$ but
archives $F_\mathrm{anchor}$.
Insufficient adjustment reflects incomplete modification of the
selection rule: $\mathcal{S}$ continues to assign partial weight
to $F_\mathrm{anchor}$ under conditions similar to the original
anchoring context.

\subsection{\textit{C. elegans}: Minimal Stack Depth}

The predict-and-correct chemotaxis loop of \textit{C.~elegans}
\cite{bargmann2006} may implement a minimal stack of depth 2:
the current directional hypothesis and the most recent prior.
The selection rule is implemented in the gain dynamics of the
gradient detection circuit: when current evidence is strong,
the top frame is selected; when evidence is weak or contradictory,
the prior frame may contribute.

This would explain why \textit{C.~elegans} exhibits partial
reversion to prior trajectories under gradient reversal
conditions---a prediction of the accumulation model that
distinguishes it from the pure overwrite model.

% ============================================================
\section{Implications}
% ============================================================

\subsection{For Cognitive Science}

The accumulation model predicts that:
\begin{enumerate}
  \item Debiasing interventions that provide only corrective
  information will show incomplete and context-sensitive effects.
  \item The magnitude of bias reactivation will correlate with
  the similarity between current environmental conditions and
  the original acquisition context $C_\mathrm{err}$.
  \item High cognitive load, time pressure, and emotional arousal
  will systematically increase reactivation probability by
  degrading selection rule application.
\end{enumerate}

\subsection{For AI Systems}

In artificial systems, the accumulation model predicts:
\begin{enumerate}
  \item Fine-tuning that installs a corrective frame $F^*$ does
  not eliminate the prior frame $F_\mathrm{err}$ from the weight
  space; it modifies the selection dynamics.
  \item Under distribution shift conditions that resemble the
  original training distribution, the prior frame may be
  preferentially activated.
  \item Reliable debiasing requires not only corrective training
  but modification of the selection mechanism---e.g., via
  attention reweighting or gating mechanisms that suppress
  prior-context activation.
\end{enumerate}

\subsection{Connection to $d_\mathrm{cog}$}

In the terms of \cite{minamoto2026topology}, the layered
interpretive stack implies that the grounding-induced topology
$\tau_A$ of an agent is not a single topology but a \emph{mixture}:
a weighted combination of the topologies induced by each frame
in the stack, with weights determined by the selection rule.
Under the persistence-signature formulation of
\cite[Def.~\ref{def:dcog-signature}]{minamoto2026topology}, the
cognitive distance $d_\mathrm{cog}(A, B)$ between agents with
different stack histories decomposes into components that track
distinct aspects of this mixture: the $d_\mathrm{GH}$ component
reflects worst-case distortion of the currently-selected frame,
while the $d_\mathrm{bn}^{(k)}$ components reflect persistence-level
alignment that integrates over the active stack region. Differences
arising from divergent stack histories therefore show up as
$d_\mathrm{bn}^{(k)}$ differences that would be invisible to a
scalar GH measurement --- a concrete point of contact between the
accumulation model of this paper and the measurement framework of
Paper 1.

More explicitly, a companion experiment reported in
\cite[App.~\ref{app:overwrite-h1}]{minamoto2026overwrite} exhibits
a closed-form relationship between the geometry of an agent's
interpretive-frame $H_1$ cycle and its persistence lifetime: under
a constructed metric setting, the lifetime is
$1 - \sqrt{1 - \cos g(\Theta)}$, where $g$ is the maximal adjacent
gap in the cycle. If the accumulation model is correct, stack
histories that install new frames on top of old ones modify
$g(\Theta)$ through the position of the newly-installed concept
in the ring structure of the previously-dominant frame. Whether
the new frame overwrites a critical node or a redundantly-placed
one determines whether the prior frame's $H_1$ feature survives or
collapses. This gives the accumulation/overwrite distinction a
numerical signature that neither paper could supply in isolation
from the measurement framework of Paper 1.

The persistence signature is therefore not a dynamic extension of
$d_\mathrm{cog}$ in the strong sense of being time-indexed; it is
a \emph{structurally richer static snapshot} whose variation across
stack-history-differing agents captures the accumulation-specific
phenomena this paper addresses. A fully time-indexed extension
(e.g.\ via zigzag persistence) is left to future work.

% ============================================================
\section{Pilot Empirical Evidence}
\label{sec:pilot}
% ============================================================

The accumulation model predicts that prior representational structure
persists after training rather than being deleted.
We report a pilot study providing direct structural evidence for
this claim in language models.

\paragraph{Setup.}
We applied RSA to Pythia-70m \cite{biderman2023} at three training
stages: step~0 (random initialization), step~1000 (early training),
step~143000 (reference).
Pairwise concept dissimilarity matrices were computed over
$n = 29$ probe concepts at hidden layers $\ell \in \{2, 4, 6\}$,
averaged over three prompt templates.
Structural alignment was measured via second-order Spearman
correlation $\rho_\mathrm{cog}$.

\paragraph{Results.}
Three findings bear directly on the accumulation model.

\smallskip
\noindent\textbf{Finding 1: Monotone structural increase.}
$\rho_\mathrm{cog}$ increases monotonically at all layers
(L6: step-0 $= 0.017$; step-1000 $= 0.535$).
If training were pure overwrite, step-0 and step-1000 would share
no structural similarity with the reference beyond noise.
The monotone increase is consistent with the accumulation model:
prior structure is retained and built upon.

\smallskip
\noindent\textbf{Finding 2: Non-zero baseline.}
Step-0 (random initialization) shows positive but weak correlation
with the trained reference at all layers.
Under the accumulation model, this reflects residual structural
contribution of pre-training initialization to the final stack.
Under the correction model, no such similarity is expected.

\smallskip
\noindent\textbf{Finding 3: Cross-architecture convergence.}
$\rho_\mathrm{cog}(\text{GPT-2}, \text{Pythia-trained})
\approx 0.42$ versus
$\rho_\mathrm{cog}(\text{GPT-2}, \text{Pythia-step-0})
\approx 0.10$ across all layers.
Two architectures trained on different corpora converge to similar
relational geometry, consistent with the accumulation model's
prediction that training installs distribution-determined frames.

\paragraph{Interpretation.}
These results are consistent with the subsection on AI systems
(Section~6.2): fine-tuning modifies selection dynamics without
eliminating prior frame structure.
Finding 2 provides an initial quantitative handle on
``selection instability as the source of bias'' (Section~3):
the residual structure of step-0 in the trained reference is
precisely the accumulated prior that can be reactivated under
conditions resembling the original acquisition context.
Direct mechanistic tests (circuit persistence, context-dependent
reactivation) remain as future work.
All code is available at
\url{https://github.com/minamominamoto/grounding-without-reality}.

% ============================================================
\section{Open Problems}
% ============================================================

\begin{description}
  \item[OP1.] \textbf{Formal characterization of the selection rule.}
    What is the computational form of $\mathcal{S}$?
    Candidate: a softmax over frame-context similarity scores,
    where similarity is computed between current environmental
    state and each frame's acquisition context.

  \item[OP2.] \textbf{Stack depth in biological systems.}
    What is the effective stack depth in human cognition?
    In \textit{C.~elegans}?
    Is there a minimum depth required for flexible behavior?

  \item[OP3.] \textbf{Measurement of reactivation probability.}
    How can the selection probability of an archived frame be
    measured empirically?
    Candidate: vary environmental conditions systematically and
    measure behavioral or neural signatures of frame reactivation.

  \item[OP4.] \textbf{SSC and NAC joint optimization.}
    What is the relationship between NAC (pre-commitment suspension)
    and SSC (post-commitment suppression)?
    Is there a tradeoff: systems that maintain longer suspension
    may have shallower stacks and therefore lower reactivation risk?

  \item[OP5.] \textbf{Fine-tuning as accumulation.}
    Does fine-tuning in large language models implement accumulation
    rather than correction?
    If so, what are the signatures of prior-frame reactivation
    in fine-tuned models, and how can they be measured via
    mechanistic interpretability tools \cite{anthropic2025circuit}?
\end{description}

% ============================================================
\section{Conclusion}
% ============================================================

The correction model of abduction is wrong in a precise and
consequential way: it assumes that prior interpretive frames
are deleted upon revision.
They are not.
They are archived in a layered interpretive stack, where they
remain as latent selection risks.

What we call ``correcting'' an abductive error is more accurately
described as installing a new top-of-stack frame and modifying
the selection rule.
The error persists; what changes is its probability of being
foregrounded.

This has a direct implication for the theory of bias.
Bias is not the presence of an erroneous frame.
It is selection instability: the erroneous frame being selected
under conditions where the correct frame should dominate.
And it has a direct implication for the design of reliable
cognitive systems, biological or artificial: the relevant
capacity is not error deletion (impossible) but selection
stability under adversarial conditions (achievable in principle).

\medskip
\noindent
\textit{Abductive errors are not eliminated by correction.\\
They persist as latent frames. They can only be suppressed.}

% ============================================================
\section*{Acknowledgments}
% ============================================================

No funding was received for this work.

% ============================================================
\begin{thebibliography}{99}

\bibitem{minamoto2026topology}
Minamoto, M. (2026).
On the topology of symbol grounding: Sensory space structure as a
determinant of cognitive representation.
Preprint (not yet on arXiv).
\url{https://github.com/minamominamoto/topology-of-grounding}

\bibitem{minamoto2026cgla}
Minamoto, M. (2026).
Reconfigurable dataflow topology as a structural condition for
universal sensory grounding.
Preprint (not yet on arXiv).
\url{https://github.com/minamominamoto/topology-of-grounding}

\bibitem{minamoto2026overwrite}
Minamoto, M. (2026).
Abduction as overwrite: Negative capability as a structural
condition for reliable inference.
Preprint (not yet on arXiv).
\url{https://github.com/minamominamoto/topology-of-grounding}

\bibitem{harman1965}
Harman, G. (1965).
The inference to the best explanation.
\textit{Philosophical Review}, 74(1), 88--95.

\bibitem{ecker2022}
Ecker, U. K. H., et al. (2022).
The psychological drivers of misinformation belief and its resistance
to correction.
\textit{Nature Reviews Psychology}, 1, 13--29.

\bibitem{loftus1979}
Loftus, E. F., \& Palmer, J. C. (1974).
Reconstruction of automobile destruction: An example of the
interaction between language and memory.
\textit{Journal of Verbal Learning and Verbal Behavior}, 13(5),
585--589.

\bibitem{epley2001}
Epley, N., \& Gilovich, T. (2001).
Putting adjustment back in the anchoring and adjustment heuristic:
Differential processing of self-generated and experimenter-provided
anchors.
\textit{Psychological Science}, 12(5), 391--396.

\bibitem{bargmann2006}
Bargmann, C.~I. (2006).
Chemosensation in \textit{C.~elegans}.
\textit{WormBook}, 1--29.

\bibitem{anthropic2025circuit}
Lindsey, J., et al. (2025).
Circuit tracing: Revealing computational graphs in language models.
Anthropic Technical Report.
\url{https://transformer-circuits.pub/2025/attribution-graphs/methods.html}

\bibitem{biderman2023}
Biderman, S., et al.\ (2023).
Pythia: A suite for analyzing large language models across training and scaling.
\textit{Proceedings of ICML 2023}.

\bibitem{radford2019language}
Radford, A., et al.\ (2019).
Language models are unsupervised multitask learners.
\textit{OpenAI Blog}, 1(8).

\end{thebibliography}

\end{document}
